\subsection*{问题 1}
\paragraph{(a)}\label{sol:att-a}
对归一化分母使用商法则时，每个指数导数多出 $1/\tau$，因此
$\partial P_j/\partial S_k=P_j(\mathbf1[j=k]-P_k)/\tau$。
再应用链式法则，得到 $D_k=P_kG^\top(V_k-Y)/\tau$。对 $k$ 求和后，两项分别为 $G^\top Y/\tau$ 和 $G^\top Y\sum_kP_k/\tau$，二者相消。分数增加常数时，分子与分母乘以同一个正数，概率不变。

当 $\tau\to\infty$ 时，每个允许位置的指数趋于 $1$，概率趋于 $1/|J_i|$。当最大分数唯一且与其他分数的差至少为 $\Delta>0$ 时，非最大位置的概率为 $O(\exp(-\Delta/\tau))$。固定有限值向量且上游梯度有界时，分数梯度也趋于零，因为指数衰减快于 $1/\tau$ 的增长。若有并列最大值，对单个原始分数的 Jacobian 元素可以随 $1/\tau$ 增大，不能沿用唯一最大值的结论。

\paragraph{(b)}\label{sol:att-b}
原例的值梯度为 $P^\top G=(1/4,3/4)$。只允许第一个位置时，概率不随分数改变，因此分数梯度为零。此时上游梯度为 $-2$，值梯度为 $(-2,0)$。若损失在该输出处的导数也为零，则该行的值梯度相应为零。

\paragraph{(c)}\label{sol:att-c}
答案查询通过 $P^\top G$ 向视觉值传递梯度。视觉位置没有直接监督标签不会自动删除这一计算路径。若该视觉位置被所有有监督查询遮蔽，则其值分支在这一层的梯度为零；其他层或其他损失的路径需要分别检查。

\paragraph{(d)}\label{sol:att-d}
对允许位置减去最大分数，计算指数和归一化；其他位置返回精确的零。向量-Jacobian 乘积为 $P\odot(g-\sum_jP_jg_j)$。中心差分检验应使用有限 logits，保持掩码固定，并避免把微分一个离散掩码作为该测试的目标。测试还应检查全遮蔽输入按接口抛出异常。

\paragraph{(e)}\label{sol:att-e}
设排列矩阵为 $\Pi$，使用 $K'=\Pi K$、$V'=\Pi V$ 并同样排列掩码列。则 $S'=S\Pi^\top$，逐行归一化给出 $P'=P\Pi^\top$，所以 $P'V'=P\Pi^\top\Pi V=PV$。只重排值时，例题的输出从 $2$ 变为 $0$，因为权重仍为 $(1/4,3/4)$，值变为 $(3,-1)$。第一个测试检查一致的轴重排，第二个测试检查分数与值的对应关系。

\subsection*{问题 2}
\paragraph{(a)}\label{sol:flow-a}
条件误差的交叉项为零，最优预测等于条件均值。利用独立性得到 $q_t=(1-t)^2s_0^2+t^2s_1^2$ 和 $c_t=ts_1^2-(1-t)s_0^2$。残差 $R=U-(m_1-m_0)-(c_t/q_t)(X_t-m_t)$ 与 $X_t$ 的协方差为零。两者联合高斯，故条件残差均值为零，得到所需速度。

反例可取 $X$ 在 $\{-1,0,1\}$ 上等概率分布，$R=X^2-2/3$。由对称性，$\E[R]=0$ 且 $\Cov(R,X)=\E[X^3]-(2/3)\E[X]=0$；但是 $\E[R\mid X]=X^2-2/3$，在三个取值处均不等于零。因此，推导中的联合高斯条件承担了从不相关到独立的关键步骤。

\paragraph{(b)}\label{sol:flow-b}
$q_t'=2c_t$ 给出 $s_t'=c_t/s_t$。对仿射映射求导即可得到 $v(t,\Phi_t)$。最终映射为 $m_1+(s_1/s_0)(x_0-m_0)$，其与初值的协方差为 $s_0s_1$；独立训练端点的协方差为零。边缘分布只规定单个时刻的分布，不唯一确定两个时刻的联合分布。

\paragraph{(c)}\label{sol:flow-c}
三个时刻的方差分别为 $1,1/2,1$；速度分别为 $-x,0,x$。另一条路径可令 $X_t=X_0$，始终保持标准高斯，速度恒为零。它满足相同端点分布，采用不同的中间分布和端点构造。

\paragraph{(d)}\label{sol:flow-d}
使用解析末点 $3.3$，分别计算离散解的绝对误差。固定步数时，中点法每步计算两次速度，Euler 法计算一次。报告应同时列方法、步数、速度求值次数和误差。有限差分在内部时刻比较解析流导数与速度场；端点使用单边差分或只检查初末值。测试结果仅适用于本题明确的光滑高斯速度场。

固定参数的标量计算结果如下。此表只比较本题设置，不构成方法一般收敛阶数的证明。

\begin{center}
\begin{tabular}{lrrr}
\toprule
方法 & 步数 & 速度求值次数 & 末点绝对误差 \\
\midrule
Euler & 50 & 50 & 0.03299895 \\
Euler & 100 & 100 & 0.01660381 \\
中点法 & 50 & 100 & 0.00000145 \\
中点法 & 100 & 200 & 0.00000018 \\
\bottomrule
\end{tabular}
\end{center}
